\begin{ExerciseBlock}
\ExerciseHeading{01}{Erzwungene Schwingung}{forced_oscillations_driven_spring}
Eine sinusförmige Kraft mit Amplitude $F_0=\qty{5.5e-3}{\newton}$ regt ein Federpendel mit $m=\qty{0.15}{\kilogram}$, $D=\qty{1.10}{\newton\per\metre}$ und Abklingkonstante $\delta=\qty{0.096}{\per\second}$ an.
\begin{enumerate}
\item Berechnen Sie die Kreisfrequenz $\omega_0$ des freien ungedämpften und die Kreisfrequenz $\omega$ des freien gedämpften Systems.
\item Bestimmen Sie die Resonanzkreisfrequenz $\omega_\mathrm{res}$ und die Resonanzamplitude $x_\mathrm{res}$.
\item Wie groß ist die statische Auslenkung $x_\mathrm{stat}$?
\item Bestimmen Sie die Güte $Q$ auf zwei Arten.
\item Bestimmen Sie den Phasenverzug zwischen Auslenkung und Erregerkraft bei $\Omega=\omega_0$ und bei $\Omega=\omega_\mathrm{res}$.
\end{enumerate}
\end{ExerciseBlock}
\UseExerciseSolution{forced_oscillations_driven_spring}
